\documentclass[11pt]{article}

% Packages
\usepackage[margin=1in]{geometry}
\usepackage{amsmath, amsfonts, amssymb}

% Title
\title{Linear Algebra Mastery \\ \large Vectors and Vector Operations}
\author{Evans N. Karago}
\date{\today}

% Begin Document
\begin{document}
\maketitle
	\section{Introduction}
		A vector is an ordered list of  numbers with both magnitude and direction. Vectors are very common in machine learning applications. Whenever a linear regression model is trained, the predictions are presented using vector-linear combination method e.g. $\hat y = w_1x_1 + w_2x_2 + \cdots + w_nx_n$. There are many other applications of vectors in real world, they are not just limited to linear regression modelling.  Some of them includes:
		\begin{enumerate}
			\item Scaling a feature, done in feature normalization
			\item Understanding span of all possible predictions a model can make. 
			\item Compute model weights, from $A\vec x = b$, given that $\vec x$ is a vector with the model weights. 
		\end{enumerate}
	
	\section{Formal Definition}
	\subsection{Vector addition}
	
	\[
		\vec u + \vec v = \begin{pmatrix}
		u_1 + v_1\\
		u_2+ v_2\\
		\vdots \\
		u_n + v_n\\
\end{pmatrix} \qquad 
\text{Given that: } \vec u , \vec v \in \mathbb{R}^n			
	\]
	\textbf{Note: } Vector addition is \textit{commutative} meaning $\vec u + \vec v = \vec v + \vec u$\\[0.15in]
	\subsection{Scalar multiplication}
	\[
	c\vec v = 
	\begin{pmatrix}
		cv_1 \\
		cv_2\\
		\vdots \\
		cv_n\\
	\end{pmatrix}	\qquad
	\text{Given that: } c \in \mathbb{R}, \quad \vec v \in \mathbb{R}^n
	\]

Scalar multiplication meaning:
\begin{enumerate}
	\item If $c > 1$: Positive stretching
	\item If $0 < c < 1$: Diminishing/shrinking vector
	\item If $c < 0$: Flipping
	\item If $c = 0$: Collapsing to a zero vector
\end{enumerate}	

	\subsection{Linear combination}
	Scaling vectors and adding them together. 
	\[
		c_1\vec v_1 + c_2\vec v_2 + \cdots + c_n \vec v_n	
	\]
	\[
		\text{Where: }	c_i \in \mathbb{R} \text{ and } \vec v_i \in \mathbb{R}^n
	\]
	
	\subsection{Span}
	Span is a set of all linear combinations of a set of vectors.
	
	\[
		span\{\vec v_1 , \vec v_2\} = span \{c_1 \vec v_1 + c_2 v_2\mid c_i \in \mathbb{R}, \quad \vec v_i \in \mathbb{R}^n\}	
	\]
	
Geometry explanation:
\begin{enumerate}
	\item Span of a non zero vector is a line through the origin.
	\item Span of a two non-parallel vectors,  spans the entire $\mathbb{R}^2$ space.
	\item Span of two parallel vectors is a line in $\mathbb{R}^2$
\end{enumerate}
	\section{Python Results}
	Running the script with $\vec u = (1, 3)$ and $\vec v = (2, -1)$ produced $\vec v + \vec u  =  (3, 2)$, confirming commutative, $\vec v + \vec u = (3,2)$. The rank check confirmed that $\vec v_1 = (1,0)$ and $\vec v_2 = (0,1)$ spans $\mathbb{R}^2$. The two vectors are not parallel, so they provide new information and spans the entire plane. $\vec v_3 = (2, 4)$ and $\vec v_4 = (1,2)$ had a rank = 1, meaning that the two vectors are parallel, therefore the two vectors form a line in  $\mathbb{R}^2$.
	
	\section{ML Connection}
	\begin{enumerate}
		\item Adding features in a feature matrix 
		\item Normalizing features
		\item Presenting predictions from a linear regression model. Linear regression predictions is a linear combination of weights and feature vector e.g. $\hat y = w_1 x_1 + w_2 x_3 + \cdots + w_nx_n$
		\item Span of a feature matrix should be the same as the one for target vector. If the target vector lies outside the span, no exact solution exists.
	\end{enumerate}
	
	
	\section{Reflection}
Todays lesson reinforced my understanding of vectors as a fundamental building block for linear algebra and machine learning.  The concept I found challenging today was the formal notation of span. Specifically distinguishing between set label $span\{\vec v_1, \vec v_2\}$ and linear combinations $span\{c_1\vec v_1 + \cdot + c_n\vec v_n\}$. What fascinate me the most, is how linear regression model predictions are presented using linear combination of weights and features.
\end{document}